% TOP_PROBLEM: {"id": 30000183, "title": "Strong Nakayama Conjecture", "category_id": 4, "category": {"id": 4, "name": "algebra", "display_name": "Algebra", "description": "Group theory, ring theory, field theory, and algebraic structures.", "slug": "algebra", "order_index": 4, "created_at": "2026-07-31T15:26:25.671Z"}, "status": "open"}
% FIRST_POSTED: 2026-10-06
% ATTEMPT_STATUS: unresolved
% OPENAI_CLAIM: {"schema_version":1,"source_repo":"https://github.com/openai/math","source_commit":"adc7f1241b42e322a6451854ab7e4b4c146bf78a","release_date":"2026-10-06","match":"partial","resolution":"partial","independently_reviewed":false,"second_pass":"n/a","families":[{"family":"198","title":"A counterexample to finitistic-dimension finiteness","subject":"Algebra","manuscripts":[{"title":"An algebra of infinite little finitistic dimension","pdf_path":"preprints/An-algebra-of-infinite-little-finitistic-dimension-September-23-2026/paper.pdf","readme_path":"preprints/An-algebra-of-infinite-little-finitistic-dimension-September-23-2026/README.md","date":"2026-09-23","theorem_ref":"Theorem 1.1"}],"lean":{"doc_path":"lean/docs/198.md","comparators":["lean/ComparatorChallenges/FinitisticAsymmetry.json","lean/ComparatorChallenges/FinitisticAsymmetry.lean","lean/ComparatorChallenges/LittleFinitistic.json","lean/ComparatorChallenges/LittleFinitistic.lean"],"declarations":["OAI.LittleFinitistic.Main.exists_counterexample","OAI.LittleFinitistic.extreme_left_right_asymmetry"],"covers_main_theorem":true},"reasoning_trace":null},{"family":"199","title":"Counterexamples to Auslander–Reiten, Tachikawa and related homological conjectures","subject":"Algebra","manuscripts":[{"title":"A counterexample to Tachikawa's second conjecture","pdf_path":"preprints/A-counterexample-to-Tachikawas-second-conjecture-September-23-2026/paper.pdf","readme_path":"preprints/A-counterexample-to-Tachikawas-second-conjecture-September-23-2026/README.md","date":"2026-09-23","theorem_ref":"Theorem 1.1"}],"lean":{"doc_path":"lean/docs/199.md","comparators":["lean/ComparatorChallenges/AuslanderReiten.json","lean/ComparatorChallenges/AuslanderReiten.lean","lean/ComparatorChallenges/Tachikawa.json","lean/ComparatorChallenges/Tachikawa.lean"],"declarations":["OAI.ArExplicit.Statement.main","OAI.Tachikawa.main_theorem"],"covers_main_theorem":true},"reasoning_trace":null}]}
\documentclass[11pt]{article}
\usepackage[margin=25mm]{geometry}
\usepackage{fontspec}
\setmainfont{Latin Modern Roman}
\usepackage{amsmath,amssymb,mathtools}
\usepackage{unicode-math}
\setmathfont{Latin Modern Math}
\usepackage{enumitem}
\usepackage{xurl,hyperref,bookmark}
\hypersetup{colorlinks=true,urlcolor=blue}
\setcounter{secnumdepth}{0}
\setlength{\parindent}{0pt}
\setlength{\parskip}{5pt}
\setlength{\emergencystretch}{3em}
\begin{document}
\section*{Strong Nakayama Conjecture}

\subsection{OpenAI's claim}
The linked manuscripts present the following claims by OpenAI: A single finite-dimensional complex algebra is claimed to have finite modules whose finite projective dimensions are unbounded; A finite-dimensional symmetric algebra is claimed to carry a nonprojective finite module without positive self-extensions, refuting Tachikawa second-conjecture rigidity. Together they provide the following support within the scopes described below.

\subsection{Outline of the argument}
One fixed finitely presented algebra encodes arbitrarily long terminating processes through a twisted derived tensor functor. The subsequent finite-dimensional realization converts these tensor histories into ordinary projective resolutions, retaining unbounded termination lengths without changing the final algebra.

The construction strengthens complete-resolution vanishing on both sides of degree zero before passing to a trivial extension. A transfer calculation then produces a symmetric algebra and preserves the nonprojective module together with its positive Ext vanishing.

\subsection{Scope relative to this problem}
The finite-dimensional complex example is an Artin algebra with infinite little finitistic dimension, refuting the first clause. The separate strong-Nakayama clause quantifies over Ext in degree zero as well and is not established or refuted by this main theorem. The bundled question has two clauses. The cited corollary gives a nonzero simple S with Ext\textasciicircum{}i(S,Gamma)=0 for all i>=0, including Hom; the cited result gives modules of every finite projective dimension, hence infinite finitistic dimension. Both clauses are refuted. This entry records partial support; complete coverage of the question is not established by this catalogue.

\subsection{Formal verification}
For family 198, the Lean scope includes unbounded finite projective dimensions over one complex algebra and a separate extreme left/right finitistic and injective-generation asymmetry. The listed Comparator files are \texttt{FinitisticAsymmetry.lean}, \texttt{LittleFinitistic.lean}. Named declarations include \texttt{OAI.LittleFinitistic.Main.exists\_counterexample}, \texttt{OAI.LittleFinitistic.extreme\_left\_right\_asymmetry}. For family 199, the documented entries cover the characteristic-two rational-function-field Auslander--Reiten example, including field-extension persistence, and a separate symmetric Tachikawa counterexample. The listed Comparator files are \texttt{AuslanderReiten.lean}, \texttt{Tachikawa.lean}. Named declarations include \texttt{OAI.ArExplicit.Statement.main}, \texttt{OAI.Tachikawa.main\_theorem}. This site has not executed the Lean code. The scope descriptions do not certify the remaining manuscript arguments.

\subsection{Status caveat}
OpenAI attributes these manuscripts to work by an internal model. Its release cautions that arguments without formalization can contain mistakes. This catalogue has not independently checked the proofs; the labels report the mathematical claims attributed to OpenAI.

\subsection{Sources}
\begin{itemize}
\item OpenAI, \emph{An algebra of infinite little finitistic dimension}, 2026-09-23; citation key \texttt{OAI:An-algebra-of-infinite-little-finitistic-dimension-September-23-2026}. \href{https://github.com/openai/math/blob/adc7f1241b42e322a6451854ab7e4b4c146bf78a/preprints/An-algebra-of-infinite-little-finitistic-dimension-September-23-2026/paper.pdf}{PDF (GitHub)}; \href{https://github.com/openai/math/raw/adc7f1241b42e322a6451854ab7e4b4c146bf78a/preprints/An-algebra-of-infinite-little-finitistic-dimension-September-23-2026/paper.pdf}{PDF (direct)}; \href{https://github.com/openai/math/blob/adc7f1241b42e322a6451854ab7e4b4c146bf78a/preprints/An-algebra-of-infinite-little-finitistic-dimension-September-23-2026/README.md}{Manuscript page and citation}; \href{https://github.com/openai/math/blob/main/preprints/An-algebra-of-infinite-little-finitistic-dimension-September-23-2026/paper.pdf}{Latest version}.
\item Family 198: \href{https://github.com/openai/math/blob/adc7f1241b42e322a6451854ab7e4b4c146bf78a/lean/docs/198.md}{Lean scope}.
\item \href{https://github.com/openai/math/blob/adc7f1241b42e322a6451854ab7e4b4c146bf78a/lean/ComparatorChallenges/FinitisticAsymmetry.json}{Comparator FinitisticAsymmetry.json}.
\item \href{https://github.com/openai/math/blob/adc7f1241b42e322a6451854ab7e4b4c146bf78a/lean/ComparatorChallenges/FinitisticAsymmetry.lean}{Comparator FinitisticAsymmetry.lean}.
\item \href{https://github.com/openai/math/blob/adc7f1241b42e322a6451854ab7e4b4c146bf78a/lean/ComparatorChallenges/LittleFinitistic.json}{Comparator LittleFinitistic.json}.
\item \href{https://github.com/openai/math/blob/adc7f1241b42e322a6451854ab7e4b4c146bf78a/lean/ComparatorChallenges/LittleFinitistic.lean}{Comparator LittleFinitistic.lean}.
\item OpenAI, \emph{A counterexample to Tachikawa's second conjecture}, 2026-09-23; citation key \texttt{OAI:A-counterexample-to-Tachikawas-second-conjecture-September-23-2026}. \href{https://github.com/openai/math/blob/adc7f1241b42e322a6451854ab7e4b4c146bf78a/preprints/A-counterexample-to-Tachikawas-second-conjecture-September-23-2026/paper.pdf}{PDF (GitHub)}; \href{https://github.com/openai/math/raw/adc7f1241b42e322a6451854ab7e4b4c146bf78a/preprints/A-counterexample-to-Tachikawas-second-conjecture-September-23-2026/paper.pdf}{PDF (direct)}; \href{https://github.com/openai/math/blob/adc7f1241b42e322a6451854ab7e4b4c146bf78a/preprints/A-counterexample-to-Tachikawas-second-conjecture-September-23-2026/README.md}{Manuscript page and citation}; \href{https://github.com/openai/math/blob/main/preprints/A-counterexample-to-Tachikawas-second-conjecture-September-23-2026/paper.pdf}{Latest version}.
\item Family 199: \href{https://github.com/openai/math/blob/adc7f1241b42e322a6451854ab7e4b4c146bf78a/lean/docs/199.md}{Lean scope}.
\item \href{https://github.com/openai/math/blob/adc7f1241b42e322a6451854ab7e4b4c146bf78a/lean/ComparatorChallenges/AuslanderReiten.json}{Comparator AuslanderReiten.json}.
\item \href{https://github.com/openai/math/blob/adc7f1241b42e322a6451854ab7e4b4c146bf78a/lean/ComparatorChallenges/AuslanderReiten.lean}{Comparator AuslanderReiten.lean}.
\item \href{https://github.com/openai/math/blob/adc7f1241b42e322a6451854ab7e4b4c146bf78a/lean/ComparatorChallenges/Tachikawa.json}{Comparator Tachikawa.json}.
\item \href{https://github.com/openai/math/blob/adc7f1241b42e322a6451854ab7e4b4c146bf78a/lean/ComparatorChallenges/Tachikawa.lean}{Comparator Tachikawa.lean}.
\item \href{https://github.com/openai/math/blob/adc7f1241b42e322a6451854ab7e4b4c146bf78a/CONTENTS.md}{OpenAI catalogue entry}.
\end{itemize}
\end{document}
